Recent research suggests that increasing inference-time computation through "thinking tokens" or "pause tokens" can significantly improve the reasoning capabilities of Large Language Models (LLMs). However, most existing approaches rely on specialized training or reinforcement learning. In this work, we investigate whether these benefits can be achieved through simple prompt-level instructions—a paradigm we call \textit{Between-Sentence Thinking} (\bst). We evaluate \gpt on \tqa and complex creative writing tasks under three conditions: standard generation, meaningless filler tokens between sentences (\bstfiller), and explicit internal rationales (\bstrationale). Contrary to the "dot-by-dot" computation hypothesis, we find that prompted \bst consistently degrades performance across all benchmarks. Qualitative analysis reveals a pervasive failure mode we term \textit{Model Shirking}: when faced with the additional computational overhead of generating thinking tokens, the model minimizes its total effort by producing significantly shorter, less detailed, and less informative responses. Control models outperformed \bst variants in 95\% of \tqa cases and 100\% of creative writing tasks. Our results suggest that without proper alignment or reinforcement, prompting for extra computation triggers a "laziness" response that overrides any potential reasoning gains.